% ============================================================
\chapter{Étude et réalisation du Sprint 5}
\label{chap:sprint5}
% ============================================================

\section{Introduction}
\label{sec:intro_sprint5}

Ce chapitre présente l'étude et la réalisation du
\textbf{cinquième sprint} de la plateforme FinBank. Après le
développement des fonctionnalités bancaires et de l'assistant
conversationnel, ce sprint porte sur l'orchestration des processus
métiers et la réalisation des interfaces de la plateforme.

Les travaux concernent l'intégration du moteur \textbf{Flowable}, la
modélisation des processus BPMN de validation des virements, de
détection de fraude et de mise à jour des données. Ils comprennent
également le développement des interfaces React et TypeScript pour les
utilisateurs et les administrateurs, ainsi que leur connexion aux API
REST.

\section{Backlog du Sprint}
\label{sec:backlog_sprint5}

\begin{table}[htbp]
\centering
\caption{Backlog du Sprint 5}
\label{tab:backlog_sprint5}
\renewcommand{\arraystretch}{1.3}
\begin{tabular}{|c|p{10cm}|c|}
\hline
\textbf{ID} & \textbf{Tâches} & \textbf{Estimation} \\
\hline
S5-T1 & Configurer Flowable et intégrer les processus BPMN dans le backend & 3 jours \\
\hline
S5-T2 & Implémenter le processus de validation des virements & 4 jours \\
\hline
S5-T3 & Implémenter les processus de détection et de traitement de la fraude & 4 jours \\
\hline
S5-T4 & Implémenter le processus de validation des mises à jour de données & 3 jours \\
\hline
S5-T5 & Exposer les tâches, processus et traces Flowable pour l'administrateur & 3 jours \\
\hline
S5-T6 & Développer les interfaces utilisateur React/TypeScript & 4 jours \\
\hline
S5-T7 & Développer les interfaces d'administration et de suivi des tâches & 3 jours \\
\hline
S5-T8 & Connecter les interfaces aux API REST et gérer les états d'erreur & 3 jours \\
\hline
\multicolumn{2}{|r|}{\textbf{Total}} & \textbf{27 jours} \\
\hline
\end{tabular}
\end{table}

\section{Conception}
\label{sec:conception_sprint5}

La conception du Sprint 5 décrit la collaboration entre l'utilisateur,
le système FinBank et l'administrateur chargé de traiter les tâches
humaines. Les processus BPMN sont déclenchés par les opérations
nécessitant une validation ou une analyse complémentaire.

\subsection{Diagramme de cas d'utilisation}
\label{subsec:usecase_sprint5}

Le diagramme présente les cas d'utilisation liés à l'exécution des
processus métier et à l'utilisation des interfaces. L'utilisateur peut
soumettre une opération et suivre son état. L'administrateur peut
consulter les tâches, examiner les détails d'un processus, prendre une
décision et consulter sa trace d'exécution.

\begin{figure}[htbp]
\centering
\IfFileExists{Figures/usecase_sprint5.png}{%
  \includegraphics[width=\textwidth,height=0.78\textheight,keepaspectratio]{Figures/usecase_sprint5.png}%
}{%
  \fbox{\parbox[c][5cm][c]{0.88\textwidth}{\centering
  Espace réservé au diagramme de cas d'utilisation du Sprint 5.}}
}
\caption{Diagramme de cas d'utilisation du Sprint 5}
\label{fig:usecase_sprint5}
\end{figure}

\subsection{Diagramme de classes du Sprint 5}
\label{subsec:classes_sprint5}

Le diagramme de classes présente les principaux objets liés à
l'orchestration des workflows. Une tâche de workflow est associée à un
processus, à un type de traitement et à un état. Les entités de
transaction et de notification permettent de conserver le résultat du
traitement et d'informer les acteurs concernés.

\begin{figure}[htbp]
\centering
\IfFileExists{Figures/classes_sprint5.png}{%
  \includegraphics[width=\textwidth,height=0.82\textheight,keepaspectratio]{Figures/classes_sprint5.png}%
}{%
  \fbox{\parbox[c][5cm][c]{0.88\textwidth}{\centering
  Espace réservé au diagramme de classes du Sprint 5.}}
}
\caption{Diagramme de classes du Sprint 5}
\label{fig:classes_sprint5}
\end{figure}

\subsection{Description textuelle : Traiter une tâche de workflow}
\label{subsec:desc_workflow_sprint5}

\begin{table}[htbp]
\centering
\caption{Cas d'utilisation : Traiter une tâche de workflow}
\label{tab:uc_workflow_sprint5}
\renewcommand{\arraystretch}{1.3}
\begin{tabular}{|p{3.2cm}|p{10.5cm}|}
\hline
\textbf{Élément} & \textbf{Description} \\
\hline
Titre & Traiter une tâche de workflow \\
\hline
Acteur & Administrateur (\texttt{ADMIN}) \\
\hline
But & Examiner une tâche sensible et enregistrer une décision. \\
\hline
Précondition & Un processus est démarré et une tâche administrateur est disponible. \\
\hline
Postcondition & La décision est enregistrée et le processus poursuit ou termine son exécution. \\
\hline
Scénario principal &
1. L'administrateur consulte la liste des tâches. \newline
2. Il ouvre le détail d'une tâche. \newline
3. Il examine les données et le contexte du processus. \newline
4. Il choisit une décision autorisée. \newline
5. Le système complète la tâche. \newline
6. Le système retourne le nouvel état du processus. \\
\hline
Scénario alternatif & Si la tâche est absente, déjà traitée ou si la décision est invalide, le système refuse l'opération. \\
\hline
\end{tabular}
\end{table}

\subsection{Diagramme de classes participantes : Validation Flowable}
\label{subsec:participants_flowable_sprint5}

Ce diagramme présente les classes intervenant lorsqu'un
administrateur consulte et complète une tâche Flowable.

\begin{figure}[htbp]
\centering
\IfFileExists{Figures/participants_flowable_sprint5.png}{%
  \includegraphics[width=\textwidth,height=0.78\textheight,keepaspectratio]{Figures/participants_flowable_sprint5.png}%
}{%
  \fbox{\parbox[c][5cm][c]{0.88\textwidth}{\centering
  Espace réservé au diagramme de classes participantes Flowable.}}
}
\caption{Diagramme de classes participantes -- Validation Flowable}
\label{fig:participants_flowable_sprint5}
\end{figure}

\subsection{Diagramme de séquence : Validation d'un virement}
\label{subsec:seq_transfer_sprint5}

Le diagramme de séquence décrit le déclenchement d'un processus de
validation et le traitement de la tâche administrative lorsque
l'opération nécessite une intervention humaine.

\begin{figure}[htbp]
\centering
\IfFileExists{Figures/seq_transfer_sprint5.png}{%
  \includegraphics[width=\textwidth,height=0.82\textheight,keepaspectratio]{Figures/seq_transfer_sprint5.png}%
}{%
  \fbox{\parbox[c][5cm][c]{0.88\textwidth}{\centering
  Espace réservé au diagramme de séquence de validation d'un virement.}}
}
\caption{Diagramme de séquence -- Validation d'un virement}
\label{fig:seq_transfer_sprint5}
\end{figure}

\section{Réalisation}
\label{sec:realisation_sprint5}

\subsection{Intégration des processus BPMN}
\label{subsec:bpmn_sprint5}

Les processus BPMN sont enregistrés dans les ressources du backend et
exécutés par Flowable. Le processus de validation des virements traite
les contrôles nécessaires avant l'exécution de l'opération. Les
processus de fraude et de mise à jour des données prennent en charge
des scénarios nécessitant une analyse, une approbation ou un rejet.

Chaque processus possède ses propres tâches, événements et chemins de
sortie. Les décisions prises au cours du workflow sont utilisées pour
mettre à jour l'état de la demande et déclencher les notifications
correspondantes.

\begin{figure}[htbp]
\centering
\IfFileExists{Figures/bpmn_processes_sprint5.png}{%
  \includegraphics[width=0.95\textwidth]{Figures/bpmn_processes_sprint5.png}%
}{%
  \fbox{\parbox[c][4cm][c]{0.88\textwidth}{\centering
  Espace réservé aux diagrammes BPMN du Sprint 5.}}
}
\caption{Processus BPMN de validation et de contrôle}
\label{fig:bpmn_processes_sprint5}
\end{figure}

\subsection{Administration des tâches et des processus}
\label{subsec:admin_flowable_sprint5}

Le contrôleur d'administration Flowable permet de lister les tâches,
les instances de processus et leur état. Il permet également de
consulter le détail et la trace d'une instance, puis de compléter une
tâche avec une décision. L'accès à ces opérations est réservé à
l'administrateur.

\begin{figure}[htbp]
\centering
\IfFileExists{Figures/admin_flowable_sprint5.png}{%
  \includegraphics[width=0.9\textwidth]{Figures/admin_flowable_sprint5.png}%
}{%
  \fbox{\parbox[c][4cm][c]{0.85\textwidth}{\centering
  Espace réservé à la capture de l'administration des tâches Flowable.}}
}
\caption{Interface d'administration des tâches et processus Flowable}
\label{fig:admin_flowable_sprint5}
\end{figure}

\subsection{Interfaces React et TypeScript}
\label{subsec:interfaces_sprint5}

Les interfaces React et TypeScript fournissent les écrans nécessaires
aux utilisateurs et aux administrateurs. Les utilisateurs disposent
notamment des vues du tableau de bord, des transactions, de
l'historique, des cartes, des investissements, de l'aide et des
sessions de conversation. Les administrateurs disposent d'un espace
dédié aux tâches, aux utilisateurs, aux journaux et aux événements de
sécurité.

\begin{figure}[htbp]
\centering
\IfFileExists{Figures/interfaces_react_sprint5.png}{%
  \includegraphics[width=0.9\textwidth]{Figures/interfaces_react_sprint5.png}%
}{%
  \fbox{\parbox[c][4cm][c]{0.85\textwidth}{\centering
  Espace réservé aux captures des interfaces React et TypeScript.}}
}
\caption{Interfaces utilisateur et administrateur}
\label{fig:interfaces_react_sprint5}
\end{figure}

\subsection{Connexion aux API REST}
\label{subsec:api_front_sprint5}

Les pages frontend communiquent avec le backend à travers des modules
API spécialisés. Ces modules centralisent les appels HTTP, les
paramètres et les types TypeScript associés aux réponses. Les états de
chargement, de succès et d'erreur sont pris en compte afin de fournir
un retour clair à l'utilisateur.

\begin{figure}[htbp]
\centering
\IfFileExists{Figures/api_front_sprint5.png}{%
  \includegraphics[width=0.9\textwidth]{Figures/api_front_sprint5.png}%
}{%
  \fbox{\parbox[c][4cm][c]{0.85\textwidth}{\centering
  Espace réservé à la capture des appels API depuis les interfaces.}}
}
\caption{Connexion des interfaces aux API REST}
\label{fig:api_front_sprint5}
\end{figure}

\section{Conclusion}
\label{sec:conclusion_sprint5}

Ce chapitre a présenté l'étude et la réalisation du Sprint 5 de
FinBank. La conception a été décrite à travers les diagrammes de cas
d'utilisation, de classes, de classes participantes et de séquence.

La réalisation a permis d'intégrer les workflows BPMN avec Flowable,
de traiter les validations de virements, les analyses de fraude et les
mises à jour de données, puis de fournir les interfaces React et
TypeScript connectées aux API REST. Le chapitre suivant sera consacré
aux tests, aux corrections et à la finalisation du projet.

